\documentclass[12pt,a4paper]{article}
\usepackage[brazilian]{babel}
\usepackage{amsmath,amssymb,amsthm}
\usepackage{geometry}
\usepackage{enumitem}
\usepackage{hyperref}
\usepackage{xcolor}
\usepackage{graphicx}
\usepackage{fancyhdr}
\usepackage{tcolorbox}
\usepackage{totpages}
\usepackage{titlesec}

% Margens ajustadas
\geometry{top=3.5cm, bottom=2.5cm, left=2.0cm, right=2.0cm, headheight=2.6cm}

% Paleta de Cores Profissional
\definecolor{primary}{RGB}{10, 45, 90}       % Azul Escuro Institucional
\definecolor{secondary}{RGB}{25, 110, 180}   % Azul Destaque
\definecolor{bglight}{RGB}{245, 247, 250}    % Cinza Claro para Caixas
\definecolor{borderlight}{RGB}{210, 220, 230}% Borda Suave

% Estilização dos Títulos de Seção
\titleformat{\section}
  {\color{primary}\normalfont\Large\bfseries}
  {\thesection}{1em}{}[\titlerule]
\titleformat{\subsection}
  {\color{secondary}\normalfont\large\bfseries}
  {\thesubsection}{1em}{}

% Caixas Estilizadas com tcolorbox
\tcbuselibrary{skins, breakable}

% Caixa para Teoremas
\newtcolorbox{teoremabox}[1][]{
  colback=bglight,
  colframe=primary,
  fonttitle=\bfseries\color{white},
  title={#1},
  arc=3pt,
  boxrule=1pt,
  enhanced,
  drop shadow=black!10
}

% Caixa para Exemplos
\newtcolorbox{exemplobox}[1][]{
  colback=white,
  colframe=secondary,
  fonttitle=\bfseries\color{white},
  title={#1},
  arc=2pt,
  boxrule=0.8pt,
  breakable
}

% Caixa de Destaque / Exemplo (Corrigida e adicionada para o documento)
\newtcolorbox{destaquebox}[1][]{
  colback=bglight,
  colframe=secondary,
  sharp corners,
  boxrule=0.8pt,
  top=3mm,
  bottom=3mm,
  left=4mm,
  right=4mm,
  breakable,
  #1
}

% Configuração de Cabeçalho e Rodapé
\pagestyle{fancy}
\fancyhf{}

% Cabeçalho
\fancyhead[L]{%
  % Certifique-se de que a imagem existe na pasta indicada ou altere o caminho
  \includegraphics[height=2cm]{Imagens/logocaruaru.png}%
}
\fancyhead[R]{%
  \small\sffamily
  \textbf{Instituto Federal de Pernambuco -- Campus Caruaru} \\
  \textcolor{secondary}{Disciplina:} Matemática VI \\
  \textcolor{secondary}{Prof.:} Cleibson José da Silva
}

% Rodapé
\fancyfoot[L]{\small\sffamily\color{gray}Material Didático -- Polinômios}
\fancyfoot[R]{\small\sffamily\color{gray}Página \thepage\ de \ref{TotPages}}

% Linhas divisórias do cabeçalho e rodapé
\renewcommand{\headrulewidth}{0.8pt}
\renewcommand{\headrule}{{\color{primary}\hrule width\headwidth height \headrulewidth}}
\renewcommand{\footrulewidth}{0.4pt}

\hypersetup{
    colorlinks=true,
    linkcolor=secondary,
    urlcolor=secondary,
    citecolor=primary
}

\begin{document}

% Título Principal do Material
\begin{center}
  \vspace*{-0.3cm}
  {\LARGE\bfseries\color{primary} Expressões Algébricas e Polinômios} \\[0.2cm]
  {\small\sffamily\color{gray} Classificação, Monômios, Operações e Divisão Algébrica}
  \vspace*{0.3cm}
  \hrule height 1.5pt \color{primary}
\end{center}

\vspace{0.3cm}

\section{Expressões Algébricas e Monômios}

Em diversas situações no estudo da Matemática e das Engenharias, é conveniente denotar números reais arbitrários por meio de letras (variáveis), possibilitando manipular algebricamente quantidades mesmo antes de conhecer seus valores numéricos exatos.

Define-se como \textbf{expressão algébrica} todo resultado de um número finito de operações fundamentais (adição, subtração, multiplicação, divisão, potenciação e radiciação) efetuadas sobre variáveis e números reais, desde que tais operações sejam bem definidas em $\mathbb{R}$.

\subsection{Classificação das Expressões Algébricas}
As expressões algébricas são classificadas segundo a natureza das operações às quais as variáveis estão submetidas:
\begin{itemize}[leftmargin=*]
    \item \textbf{Racionais Inteiras:} Não apresentam variáveis no denominador nem sob radicais. Exemplo: $E(x,y) = 3x^2 - \frac{5}{2}y^3$.
    \item \textbf{Racionais Fracionárias:} Apresentam variáveis no denominador. Exemplo: $F(x) = \frac{x+1}{x-2}$, definida para $x \neq 2$.
    \item \textbf{Irracionais:} Apresentam variáveis sujeitas à operação de radiciação. Exemplo: $G(x) = \sqrt{x+1}$, definida para $x \ge -1$.
\end{itemize}

\subsection{Monômios e Operações Elementares}
Um \textbf{monômio} é uma expressão algébrica obtida pelo produto de um número real não nulo (denominado \textbf{coeficiente}) por uma sequência finita de potências de variáveis com expoentes inteiros não negativos (denominada \textbf{parte literal}).

O \textbf{grau de um monômio} corresponde à soma dos expoentes de todas as variáveis contidas na sua parte literal. Dois ou mais monômios são ditos \textbf{semelhantes} quando possuem rigorosamente a mesma parte literal.

\begin{destaquebox}
\textbf{Exemplo 1 -- Operações com Monômios}\\
Considere os monômios $M_1 = 5ab^2$ e $M_2 = -4a^2b^3c^2$:
\begin{enumerate}[label=\alph*)]
    \item \textbf{Multiplicação:} $(5ab^2) \cdot (-4a^2b^3c^2) = (5 \cdot (-4)) \cdot (a \cdot a^2) \cdot (b^2 \cdot b^3) \cdot c^2 = -20a^3b^5c^2$.
    \item \textbf{Divisão:} Para $m, n, p \neq 0$, $\frac{45m^8n^7p^3}{9m^3n^2p} = \left(\frac{45}{9}\right) m^{8-3} n^{7-2} p^{3-1} = 5m^5n^5p^2$.
\end{enumerate}
\end{destaquebox}

\section{Representação e Estrutura dos Polinômios}

Define-se como \textbf{polinômio} em uma variável real $x$ toda expressão algébrica dada pela soma finita de monômios da forma:
\[ P(x) = a_n x^n + a_{n-1} x^{n-1} + \dots + a_1 x + a_0 \]
em que $a_n, a_{n-1}, \dots, a_0 \in \mathbb{R}$ são os coeficientes numéricos do polinômio e $n \in \mathbb{N}$.

Se $a_n \neq 0$, o expoente $n$ define o \textbf{grau do polinômio}, denotado por $\partial P$ ou $\operatorname{gr}(P)$. Quando o coeficiente do termo de maior grau for unitário ($a_n = 1$), o polinômio é dito \textbf{mônico}. Se todos os coeficientes forem nulos ($a_n = \dots = a_0 = 0$), o polinômio é \textbf{identicamente nulo} ($P \equiv 0$), não se definindo grau para este caso.

Dois polinômios $P(x)$ e $Q(x)$ são ditos \textbf{idênticos} ($P \equiv Q$) se, e somente se, os coeficientes de seus termos correspondentes de mesmo grau forem idênticos. Um número real $k$ é denominado \textbf{raiz} ou zero de $P(x)$ se $P(k) = 0$.

\begin{destaquebox}
\textbf{Exemplo 2 -- Raízes de um Polinômio}\\
Os números reais $x = 2$ e $x = 5$ são raízes do polinômio $P(x) = x^2 - 7x + 10$, pois:
\[ P(2) = 2^2 - 7(2) + 10 = 4 - 14 + 10 = 0 \]
\[ P(5) = 5^2 - 7(5) + 10 = 25 - 35 + 10 = 0 \]
\end{destaquebox}

\section{Operações Fundamentais com Polinômios}

\subsection{Adição, Subtração e Multiplicação}
\begin{itemize}[leftmargin=*]
    \item \textbf{Adição e Subtração:} Efetuam-se adicionando ou subtraindo os coeficientes dos termos de mesmo grau. O grau do polinômio resultante satisfaz a desigualdade $\partial(P \pm Q) \le \max\{\partial P, \partial Q\}$.
    \item \textbf{Multiplicação:} Multiplica-se cada termo de $P(x)$ por todos os termos de $Q(x)$, reduzindo os monômios semelhantes. O grau do produto satisfaz a igualdade $\partial(P \cdot Q) = \partial P + \partial Q$.
\end{itemize}

\begin{destaquebox}
\textbf{Exemplo 3 -- Multiplicação de Polinômios}\\
Efetue o produto de $P(x) = x^2 - 2x + 4$ por $Q(x) = -2x + 2$:
\[ P(x) \cdot Q(x) = (x^2 - 2x + 4)(-2x + 2) \]
\[ = x^2(-2x + 2) - 2x(-2x + 2) + 4(-2x + 2) \]
\[ = -2x^3 + 2x^2 + 4x^2 - 4x - 8x + 8 = -2x^3 + 6x^2 - 12x + 8 \]
\end{destaquebox}

\section{Divisão de Polinômios e Algoritmo da Divisão}

A divisão entre dois polinômios fundamenta-se na existência e unicidade de dois polinômios característicos, conforme enunciado no teorema a seguir.

\paragraph{Teorema 1 -- Algoritmo da Divisão (Euclidiana)}
Dados os polinômios $A(x)$ (dividendo) e $B(x) \neq 0$ (divisor), existem únicos polinômios $Q(x)$ (quociente) e $R(x)$ (resto) tais que:
\[ A(x) = Q(x) \cdot B(x) + R(x) \]
satisfazendo a condição de parada: $\partial R < \partial B$ ou $R(x) \equiv 0$.

\subsection{Método da Chave}
O procedimento prático para determinar $Q(x)$ e $R(x)$ consiste em:
\begin{enumerate}[label=\arabic*.]
    \item Ordenar dividendo e divisor em ordem decrescente segundo as potências de $x$.
    \item Preencher com coeficientes nulos os termos ausentes do dividendo.
    \item Dividir o termo de maior grau do dividendo pelo termo de maior grau do divisor, obtendo um termo de $Q(x)$.
    \item Multiplicar esse termo pelo divisor, inverter os sinais dos resultados e subtrair do dividendo.
    \item Repetir o ciclo até obter um resto parcial cujo grau seja estritamente menor que $\partial B$.
\end{enumerate}

\begin{destaquebox}
\textbf{Exemplo 4 -- Aplicação do Método da Chave}\\
Determine o quociente e o resto da divisão de $A(x) = 5x^4 - 4x^2 + 3x - 1$ por $B(x) = x^2 - 3x + 4$.

\textbf{Solução:} Organizando os termos no dispositivo da chave e preenchendo o termo faltante $0x^3$:

\vspace{0.3cm}
\begin{center}
\begin{tabular}{r@{\quad}l@{\quad}l}
$5x^4 + 0x^3 - 4x^2 + 3x - 1$ & \multicolumn{1}{|l}{$x^2 - 3x + 4$} & \\ \cline{2-2}
\underline{$-(5x^4 - 15x^3 + 20x^2)$} & $5x^2 + 15x + 21$ & $\leftarrow Q(x) \text{ (Quociente)}$ \\
$15x^3 - 24x^2 + 3x - 1$ & & \\
\underline{$-(15x^3 - 45x^2 + 60x)$} & & \\
$21x^2 - 57x - 1$ & & \\
\underline{$-(21x^2 - 63x + 84)$} & & \\
$\mathbf{6x - 85}$ & $\leftarrow R(x) \text{ (Resto)}$ & 
\end{tabular}
\end{center}
\vspace{0.3cm}

Portanto, o quociente obtido é $Q(x) = 5x^2 + 15x + 21$ e o resto é $R(x) = 6x - 85$. Como $\partial R = 1 < 2 = \partial B$, o algoritmo é encerrado.
\end{destaquebox}

\section*{Referências Bibliográficas}
\begin{enumerate}[label={[\arabic*]},leftmargin=*]
    \item PARENTE, Ulisses Lima; NETO, Antonio Caminha M. \textbf{Módulo de Expressões Algébricas e Polinômios -- Parte 1}. Portal da Matemática OBMEP.
    \item LIMA, Elon Lages et al. \textbf{A Matemática do Ensino Médio: Volume 3}. Coleção do Professor de Matemática. Rio de Janeiro: Sociedade Brasileira de Matemática (SBM), 2012.
\end{enumerate}

\end{document}